\documentclass[10pt,a4paper]{amsart}
\usepackage[margin=19mm]{geometry}
\usepackage{amsmath,amssymb,mathtools}
\usepackage[scr=rsfs,cal=euler]{mathalfa}
\usepackage{xfrac}
\usepackage[unicode,hidelinks,bookmarks=false]{hyperref}
\newcommand{\ie}{i.e.\ }
\newcommand{\cf}{cf.\ }
\newcommand{\resp}{respectively\ }
\newcommand{\Subsection}{\subsection}
\providecommand{\nobreakdash}{\nobreak\hbox{-}}
\setlength{\emergencystretch}{2em}
\multlinegap0pt
\usepackage{bm}
\usepackage{bbm}
\usepackage{dsfont}
\usepackage{xspace}
\usepackage{cancel}
\usepackage{mathtools}
\usepackage{amsthm}
\makeatletter
\@ifundefined{theorem}{\newtheorem{theorem}{Theorem}}{}
\@ifundefined{lemma}{\newtheorem{lemma}[theorem]{Lemma}}{}
\makeatother
\title[Finite rank perturbation of non-Hermitian random matrices: h]{Finite rank perturbation of non-Hermitian random matrices: heavy tail and sparse regimes}
\author{Yi Han}
\date{Source: arXiv:2407.21543v2 (CC BY 4.0)}
\begin{document}
\maketitle

\noindent\emph{Source and licence.} Finite rank perturbation of non-Hermitian random matrices: heavy tail and sparse regimes. Version arXiv:2407.21543v2 (CC BY 4.0), distributed under the Creative Commons Attribution 4.0 International licence, as recorded in the arXiv licence metadata for this exact version and shown on the abstract page https://arxiv.org/abs/2407.21543v2. Full rights notice: licenses/arxiv-2407.21543-CC-BY-4.0.txt.\par\smallskip

\noindent\textit{Editorial scope.} Counted unit: the Lemma whose source label is \texttt{fuckgreat} in the article (source0.tex lines 711--716, source label \texttt{fuckgreat}), delivered together with the article's own complete original proof of it (source0.tex lines 718--737, one \texttt{proof} environment of 20 lines). No mathematics is written, repaired or re-derived: statement and proof are the article's own text line for line. Required context retained uncounted: none. Every definition, display and label of those lines is retained unchanged. Omitted from this package and not redistributed: 1--710, 717--717, 738--943. Nothing inside a retained excerpt is omitted or altered: every retained body is delivered exactly as the article's lines are written, so no omission receipt of any kind accompanies this package. Citations: the retained excerpts carry no citation command at all, so no bibliography entry of the article is transcribed and no reference is redistributed. Reference adaptation: none was needed and none was made; every reference inside the delivered unit resolves inside it, so no unresolved reference or citation is produced. Printed slips kept literal: no printed word, symbol, hypothesis or number of the counted statement or of its proof was altered. Layout and class: the article's own class is \texttt{amsart} with packages including amsmath, amsthm, amscd, amsfonts, amssymb, graphicx, color, mathrsfs, hyperref, xy, slashed; the wrapper is the lane's pinned \texttt{amsart} editorial wrapper at 10pt on A4, which restates the theorem-like environments the retained text needs and adds the article's own macro definitions that the retained text needs (none), with extra wrapper packages (bm, bbm, dsfont, xspace, cancel, mathtools). The delivered numbering is the unit's own. Assets: no retained span contains a figure environment, a graphics-inclusion command, a PDF-page inclusion, a table or a TikZ picture; no image file is copied into this package, the package declares no asset dependency and no image file is redistributed with it. Licence: version arXiv:2407.21543v2 of this article is distributed under the Creative Commons Attribution 4.0 International licence, as recorded in the arXiv licence metadata for this exact version and shown on the abstract page https://arxiv.org/abs/2407.21543v2. A full rights notice is delivered at licenses/arxiv-2407.21543-CC-BY-4.0.txt. Characters: every retained excerpt is pure ASCII, so the delivered engine, XeLaTeX with its default Latin Modern text font, reports no missing character.
\begin{lemma}\label{fuckgreat}
    We have the following convergence in law:
\begin{equation}\label{eq2.14}
   \frac{\operatorname{Tr}((A_n+\sqrt{n}C_n)^k)-\operatorname{Tr}((A_n)^k)-\operatorname{Tr}((\sqrt{n}C_n)^k)}{n^{k/2}}\overset{\text{law}}{\underset{n\to\infty}\longrightarrow}0.
\end{equation}
\end{lemma}

\begin{proof} We first show that the expectation of the left hand side of \eqref{eq2.14} converges to zero. Recall that a typical term in the numerator of the LHS of \eqref{eq2.14} is of the form $w_{\mathbf{i}}=w_{i_1i_2}\cdots w_{i_ki_1}$, then in order for $\mathbb{E}[w_\mathbf{i}]\neq 0$, we need that each edge in $w_\mathbf{i}$ that are traversed by some $a_{ij}$ must be traversed twice (those travelled three times or more are negligible). Thus we can decompose the cycle $(i_1,i_2,\cdots,i_{k},i_1)$ into (a) a double cycle consisting of all the $a_{ij}$'s and some $\frac{\mu}{\sqrt{n}}$s, and (b) some other cycles attached to this cycle that consist entirely of $\frac{\mu}{\sqrt{n}}$ terms. Suppose the double cycles in (a) consist of $k_1$ vertices (and hence $k_1$ edges), and among the $k_1$ edges, $k_2$ of them are traversed by some $a_{ij}$'s and the other by $\frac{\mu}{\sqrt{n}}$.  Then there are remaining $k-2k_1$ edges to be added afterwards. This would account for an addition of strictly less than $k-2k_1$ new vertices. Thus the overall contribution to the expectation of the left hand side of \eqref{eq2.14} is, in the case $k>2k_1$: 
$$
n^{-\frac{k}{2}}\times O(n^{k_1}n^{k-2k_1-1}n^{-\frac{k-2k_1}{2}}n^{-(k_1-k_2)})=O(n^{-1}),
$$
and in the case $k=2k_1$, one must have $k_1>k_2$, so that the contribution is
$$
n^{-\frac{k}{2}}\times O(n^{k_1}n^{-(k_1-k_2)})=O(n^{-1}).
$$
This proves the expectation of the left hand side of \eqref{eq2.14} converges to $0$.

We then show the variance of the left hand side of \eqref{eq2.14} converges to $0$ as well. For two typical terms $w_{\mathbf{i}}=w_{i_1i_2}\cdots w_{i_ki_1}$ and $r_{\mathbf{j}}=r_{j_1j_2}\cdots r_{j_kj_1}$ to have non-zero covariance, we consider the graph formed by $\mathbf{i}\cup\mathbf{j}$, whose vertices and edges are the union of those from $\mathbf{i}$ and $\mathbf{j}$. Then each edge traversed by some $a_{ij}$ in $\mathbf{i}\cup\mathbf{j}$ should be traversed twice (those travelled more times are negligible), and such that there are equal number of terms in $w_\mathbf{i},r_\mathbf{j}$ travelled by $a_{ij}$'s, and they have the same subscript. Therefore $\mathbf{i}\cup\mathbf{j}$ is made up of (a) a double cycle with $k_1$ edges (each travelled twice) and $k_1$ vertices, of which $k_2$ are travelled by the $a_{ij}$'s; and (b) $2k-2k_1$ outlying edges, giving rise to at most $2k-2k_1-1$ new vertices. Then the overall contribution to the variance of the LHS of \eqref{eq2.14} is at most 
$$
n^{-k}\times O(n^{k_1}n^{2k-2k_1-1}n^{-\frac{2k-2k_1}{2}}n^{-(k_1-k_2)})=O(n^{-1}),
$$ in the case when $k<k_1$,
and 
$$
n^{-k}\times O(n^kn^{-(k_1-k_2)})=O(n^{-1}),$$
in the case when $k=k_1$ so that $k_1>k_2$,
completing the proof.
\end{proof}

\end{document}
